\section{Einleitung}
\label{sec:intro}

Technische Dokumentation entsteht heute in verteilten Teams, häufig aus
mehreren Quelltexten, Bildern und Datenquellen. Dieses Demoprojekt zeigt, wie
sich solche Bestandteile in einem einzigen \LaTeX-Projekt zusammenführen lässt
und wie \emph{LaTeX Studio} dabei unterstützt~\cite{lamport1994}.

\begin{hinweis}{Gliederungsprinzip}
  Jedes Kapitel liegt in einer eigenen Datei unter \texttt{chapters/}. Die
  Hauptdatei \texttt{main.tex} bindet sie per \texttt{\textbackslash input}
  ein. Die Gliederungs-Outline links zeigt die Struktur in Echtzeit.
\end{hinweis}

\subsection{Ziele}
\label{sec:ziele}
\begin{enumerate}[leftmargin=*, label=\alph*)]
  \item Reproduzierbare Build-Kette ohne manuelle Eingriffe.
  \item Konsistente Verweise auf Abschnitte, Abbildungen und Tabellen.
  \item Frühzeitige Fehlererkennung durch Auswertung des \LaTeX-Logs.
\end{enumerate}

\subsection{Leserleitfaden}
Abschnitt~\ref{sec:methodik} beschreibt das Vorgehen,
Abschnitt~\ref{sec:ergebnisse} wertet die Messungen aus,
Abschnitt~\ref{sec:fazit} fasst die Ergebnisse zusammen. Eine ausführliche
Diskussion der Ergebnisse findet sich in \cref{sec:diskussion}.
